

% ============================================================

% 第一章：Git 与 GitHub

% 文件：chapters/chapter01_git.tex

%

% 本章学习目标：

% 1. 理解 Git 与 GitHub 的区别

% 2. 掌握 Git 安装、配置及仓库克隆

% 3. 掌握 add、commit、push、pull 等核心命令

% 4. 学会查看修改、版本恢复及忽略文件

% 5. 为后续 Python 与机器学习项目建立版本管理基础

%

% 注意：

% 本文件是章节文件，由 main.tex 通过 \include{} 引入。

% 不要在此添加 documentclass 或 document 环境。

% ============================================================





% ============================================================

% 1.1 Git 与 GitHub 基础概念

% ============================================================



\chapter{Git 与 GitHub}

\label{chap:git_github}



\section{Git 基础概念}



\subsection{什么是 Git}



Git 是一种分布式版本控制系统

（Distributed Version Control System），

主要用于记录文件的修改历史、管理不同版本，

以及支持多人协作开发。



在 Python 编程与机器学习研究中，Git 可以用于：



\begin{itemize}

    \item 保存 Python 程序的历史版本；

    \item 比较不同算法实现之间的代码差异；

    \item 在程序修改出错时恢复稳定版本；

    \item 管理实验脚本、配置文件和研究代码；

    \item 与 GitHub、Gitee 等平台进行远程同步。

\end{itemize}



\subsection{Git 与 GitHub 的区别}



Git 是本地版本控制工具，

GitHub 是基于 Git 的远程代码托管平台。



\begin{center}

\begin{tabular}{ll}

\hline

工具 & 主要功能 \\

\hline

Git & 本地代码版本管理 \\

GitHub & 远程仓库托管与团队协作 \\

Gitee & 远程仓库托管与团队协作 \\

\hline

\end{tabular}

\end{center}



\subsection{Git 的基本工作流程}



Git 的工作流程通常包含四个部分：



\begin{enumerate}

    \item 工作区（Working Directory）；

    \item 暂存区（Staging Area）；

    \item 本地仓库（Local Repository）；

    \item 远程仓库（Remote Repository）。

\end{enumerate}



其核心命令为：



\begin{lstlisting}

git add

git commit

git push

\end{lstlisting}



% 命令解释：

% git add    ：将修改加入暂存区

% git commit ：将暂存区修改保存为本地提交

% git push   ：将本地提交上传至远程仓库



需要注意，\texttt{git commit} 仅创建本地提交，

并不意味着代码已经上传到 GitHub。





% ============================================================

% 1.2 Git 安装与配置

% ============================================================



\section{Git 安装与配置}



\subsection{检查 Git 是否安装}



在 Git Bash 或 VS Code 终端输入：



\begin{lstlisting}

git --version

\end{lstlisting}



作用：查看 Git 的安装版本。



如果终端能够正常显示版本号，

说明 Git 已安装并可以使用。



\subsection{配置用户名}



\begin{lstlisting}

git config --global user.name "Your Name"

\end{lstlisting}



% --global 表示用户级配置

% user.name 是提交记录中的作者名称

% 该名称不必与 GitHub 用户名完全一致



\subsection{配置邮箱}



\begin{lstlisting}

git config --global user.email "your-email@example.com"

\end{lstlisting}



% user.email 用于标识 Git 提交作者

% 建议使用 GitHub 账户关联的邮箱

% 也可以使用 GitHub 提供的隐私保护邮箱



\subsection{查看 Git 配置}



查看用户名：



\begin{lstlisting}

git config --global user.name

\end{lstlisting}



查看邮箱：



\begin{lstlisting}

git config --global user.email

\end{lstlisting}



查看全部有效配置：



\begin{lstlisting}

git config --list

\end{lstlisting}



\subsection{设置默认分支}



\begin{lstlisting}

git config --global init.defaultBranch main

\end{lstlisting}



% main 是常见的默认分支名称

% 该设置仅影响后续 git init 创建的新仓库

% 不会自动修改已有仓库的分支名称





% ============================================================

% 1.3 GitHub 仓库创建与克隆

% ============================================================



\section{GitHub 仓库管理}



\subsection{创建 GitHub 仓库}



在 GitHub 网站中选择 New Repository，

填写仓库名称及相关信息。



常见选项包括：



\begin{itemize}

    \item Repository name：仓库名称；

    \item Description：仓库描述；

    \item Public：公开仓库；

    \item Private：私有仓库；

    \item README：项目说明文件。

\end{itemize}



\subsection{切换工作目录}



在 Git Bash 中切换到 D 盘指定目录：



\begin{lstlisting}

cd /d/GitHub

\end{lstlisting}



查看当前工作目录：



\begin{lstlisting}

pwd

\end{lstlisting}



查看目录中的文件：



\begin{lstlisting}

ls

\end{lstlisting}



% Git Bash 使用类似 Linux 的路径格式

% Windows 的 D:\GitHub 对应 /d/GitHub



\subsection{克隆远程仓库}



命令格式：



\begin{lstlisting}

git clone <repository-url>

\end{lstlisting}



实际示例：



\begin{lstlisting}

git clone https://github.com/Zz-wutong/my-python-learning.git

\end{lstlisting}



作用：将远程仓库下载到本地，

并建立与 GitHub 仓库的关联。



\subsection{使用临时代理克隆}



如果 GitHub 网络连接失败，

且本机已经配置了 HTTP 代理，

可以使用：



\begin{lstlisting}

git -c http.proxy=http://127.0.0.1:7897 clone https://github.com/Zz-wutong/my-python-learning.git

\end{lstlisting}



% -c 表示临时设置 Git 参数

% http.proxy 指定本次命令使用的代理

% 127.0.0.1 表示本机

% 7897 是示例代理端口，应根据实际情况调整

% 该命令不会修改全局 Git 代理配置



\subsection{进入本地仓库}



\begin{lstlisting}

cd my-python-learning

\end{lstlisting}



\subsection{查看远程仓库地址}



\begin{lstlisting}

git remote -v

\end{lstlisting}



% origin 通常是克隆仓库时自动创建的远程名称

% -v 表示显示远程仓库的详细地址





% ============================================================

% 1.4 Git 本地版本管理

% ============================================================



\section{Git 本地版本管理}



\subsection{查看仓库状态}



\begin{lstlisting}

git status

\end{lstlisting}



该命令用于检查当前分支、

已修改文件、暂存文件和未跟踪文件。



常见状态说明：



\begin{itemize}

    \item \texttt{modified}：文件已被修改；

    \item \texttt{untracked}：文件尚未被 Git 跟踪；

    \item \texttt{Changes to be committed}：

          文件已经加入暂存区；

    \item \texttt{working tree clean}：

          工作区没有待提交修改。

\end{itemize}



\subsection{添加指定文件到暂存区}



\begin{lstlisting}

git add code/hello.py

\end{lstlisting}



作用：将指定文件的当前修改加入暂存区。



% 推荐初学阶段优先使用指定文件路径

% 避免意外提交其他文件



\subsection{添加当前目录的全部修改}



\begin{lstlisting}

git add .

\end{lstlisting}



作用：暂存当前目录及子目录中的新增、

修改和删除记录。



\textbf{注意：}

使用该命令前，应检查是否存在

不需要提交的配置文件或实验数据。



\subsection{取消暂存}



\begin{lstlisting}

git restore --staged code/hello.py

\end{lstlisting}



作用：将指定文件从暂存区移出，

但保留工作区中的修改。



\subsection{查看未暂存修改}



\begin{lstlisting}

git diff

\end{lstlisting}



作用：比较工作区与暂存区之间的差异。



\subsection{查看已暂存修改}



\begin{lstlisting}

git diff --staged

\end{lstlisting}



作用：检查下一次提交将包含哪些修改。





% ============================================================

% 1.5 Git Commit：创建本地提交

% ============================================================



\section{Git Commit：创建本地提交}



\subsection{创建提交}



\begin{lstlisting}

git commit -m "Update Python example"

\end{lstlisting}



% git commit 创建本地版本记录

% -m 表示直接指定提交说明

% 引号内是本次提交的描述



例如，新增试件体积计算代码后：



\begin{lstlisting}

git commit -m "Add specimen volume calculation"

\end{lstlisting}



\subsection{查看提交历史}



\begin{lstlisting}

git log

\end{lstlisting}



作用：查看提交时间、作者、Commit ID

及提交说明。



\subsection{简洁显示提交历史}



\begin{lstlisting}

git log --oneline

\end{lstlisting}



% --oneline 将每个提交压缩显示为一行

% 适合快速浏览版本历史



\subsection{查看某次提交的详细内容}



\begin{lstlisting}

git show <commit-id>

\end{lstlisting}



作用：查看指定提交引入的代码修改。





% ============================================================

% 1.6 Git Push：上传代码

% ============================================================



\section{Git Push：上传代码}



\subsection{普通上传}



\begin{lstlisting}

git push origin main

\end{lstlisting}



其中：



\begin{itemize}

    \item \texttt{push}：上传本地提交；

    \item \texttt{origin}：远程仓库名称；

    \item \texttt{main}：目标分支。

\end{itemize}



\subsection{使用临时代理上传}



\begin{lstlisting}

git -c http.proxy=http://127.0.0.1:7897 push origin main

\end{lstlisting}



% 此代理只对当前 push 命令有效

% 不改变系统环境变量

% 不修改 Git 全局代理设置



\subsection{检查上传结果}



上传完成后，可以登录 GitHub，

进入对应仓库检查文件内容和提交记录。



也可以再次执行：



\begin{lstlisting}

git status

\end{lstlisting}



注意：工作区干净不一定代表

本地提交已经成功上传，

仍需确认远程分支同步状态。





% ============================================================

% 1.7 Git Pull：获取远程更新

% ============================================================



\section{Git Pull：获取远程更新}



\subsection{拉取远程更新}



\begin{lstlisting}

git pull origin main

\end{lstlisting}



作用：获取远程分支的新提交，

并尝试整合到当前本地分支。



\subsection{使用临时代理拉取}



\begin{lstlisting}

git -c http.proxy=http://127.0.0.1:7897 pull origin main

\end{lstlisting}



\subsection{查看分支}



查看本地分支：



\begin{lstlisting}

git branch

\end{lstlisting}



查看远程分支：



\begin{lstlisting}

git branch -r

\end{lstlisting}



查看全部分支：



\begin{lstlisting}

git branch -a

\end{lstlisting}





% ============================================================

% 1.8 Git 版本恢复

% ============================================================



\section{Git 版本恢复}



\subsection{恢复未暂存的修改}



\begin{lstlisting}

git restore code/hello.py

\end{lstlisting}



作用：将工作区中的文件恢复为暂存区版本。



\textbf{警告：}

该操作会丢弃文件中尚未暂存的修改，

使用前应确认不再需要这些内容。



\subsection{撤销历史提交}



\begin{lstlisting}

git revert <commit-id>

\end{lstlisting}



作用：创建一个新的提交，

用于撤销指定提交带来的修改。



对于已经上传到 GitHub 的历史提交，

通常优先考虑使用 \texttt{git revert}。





% ============================================================

% 1.9 .gitignore 文件

% ============================================================



\section{Git 忽略文件}



\subsection{.gitignore 的作用}



\texttt{.gitignore} 用于指定不需要 Git

跟踪的文件或文件夹。



例如，Python 项目可以忽略：



\begin{lstlisting}

__pycache__/

*.pyc

.ipynb_checkpoints/

.venv/

\end{lstlisting}



如果不希望提交 VS Code 的本地配置，

还可以添加：



\begin{lstlisting}

.vscode/

\end{lstlisting}



% .gitignore 不会自动停止跟踪已经提交的文件

% 已跟踪文件需要另行从 Git 索引中移除





% ============================================================

% 1.10 完整操作实例

% ============================================================



\section{Git 完整操作实例}



假设在 VS Code 中修改了 Python 文件：



\texttt{code/hello.py}



例如新增：



\begin{lstlisting}[language=Python]

length = 400

width = 50

thickness = 0.5



volume = length * width * thickness



print("Specimen volume:", volume, "mm^3")

\end{lstlisting}



对应的 Git 操作如下：



\begin{lstlisting}

# Step 1: Check status

git status



# Step 2: Stage changes

git add code/hello.py



# Step 3: Review staged changes

git diff --staged



# Step 4: Commit locally

git commit -m "Add volume calculation"



# Step 5: Push to GitHub

git push origin main

\end{lstlisting}



需要临时代理时：



\begin{lstlisting}

git -c http.proxy=http://127.0.0.1:7897 push origin main

\end{lstlisting}





% ============================================================

% 1.11 Git 常用命令速查

% ============================================================



\section{Git 常用命令速查}



\begin{center}

\small

\begin{tabular}{ll}

\hline

命令 & 功能 \\

\hline

\texttt{git --version} & 查看版本 \\

\texttt{git status} & 查看仓库状态 \\

\texttt{git add} & 暂存修改 \\

\texttt{git commit} & 创建本地提交 \\

\texttt{git push} & 上传提交 \\

\texttt{git pull} & 拉取远程更新 \\

\texttt{git diff} & 查看修改差异 \\

\texttt{git log} & 查看提交历史 \\

\texttt{git restore} & 恢复文件 \\

\texttt{git revert} & 撤销历史提交 \\

\texttt{git remote -v} & 查看远程地址 \\

\texttt{git branch} & 查看分支 \\

\hline

\end{tabular}

\end{center}






% ============================================================
% 1.12 分支创建、切换与合并（第九课）
% ============================================================
\section{Git 分支管理}
\subsection{分支的作用}
分支（Branch）用于在不直接改变主分支的情况下开发新功能。\texttt{main} 通常保存稳定版本，\texttt{feature/...} 用于独立开发。
\subsection{创建并切换分支}
\begin{lstlisting}
git branch
# Create and switch to a feature branch
git switch -c feature/mass-calculation
# Confirm current branch
git branch
\end{lstlisting}
\subsection{在功能分支提交代码}
例如为试件添加质量计算：
\begin{lstlisting}[language=Python]
# density: g/cm^3, volume: mm^3
density = 6.45
mass = volume / 1000 * density
print("Specimen mass:", mass, "g")
\end{lstlisting}
\begin{lstlisting}
git add code/hello.py
git commit -m "Add specimen mass calculation"
\end{lstlisting}
\subsection{合并并推送主分支}
\begin{lstlisting}
git switch main
git merge feature/mass-calculation
git branch -d feature/mass-calculation
git -c http.proxy=http://127.0.0.1:7897 push origin main
\end{lstlisting}
若主分支在分支创建后没有新的独立提交，Git 通常可以执行快进合并（Fast-forward）。出现冲突时应检查冲突文件，手动解决并验证后再完成合并。

% ============================================================
% 1.13 GitHub Pull Request（第十课）
% ============================================================
\section{Pull Request 与代码审查}
\subsection{Pull Request 的作用}
Pull Request（PR）是向目标分支提出合并修改的请求，提供代码差异查看、评论、检查与审查机制。PR 与 \texttt{git pull} 不是同一个操作。
\subsection{推送功能分支}
\begin{lstlisting}
git switch main
git -c http.proxy=http://127.0.0.1:7897 pull --ff-only origin main
git switch -c feature/section-area
\end{lstlisting}
修改 \texttt{code/hello.py}：
\begin{lstlisting}[language=Python]
# Rectangular cross-sectional area in mm^2
cross_sectional_area = width * thickness
print("Cross-sectional area:", cross_sectional_area, "mm^2")
\end{lstlisting}
\begin{lstlisting}
git add code/hello.py
git commit -m "Add cross-sectional area calculation"
git -c http.proxy=http://127.0.0.1:7897 push -u origin feature/section-area
\end{lstlisting}
\subsection{GitHub 网页创建 PR}
进入仓库，点击 \textbf{Compare \& pull request}，或依次点击 \textbf{Pull requests}、\textbf{New pull request}。选择 \texttt{base: main} 和 \texttt{compare: feature/section-area}，填写标题、描述后创建 PR。
\subsection{审查与合并}
在 \textbf{Files changed} 检查代码差异，确认修改范围、公式、单位和敏感信息。审查通过后点击 \textbf{Merge pull request}，再确认合并。网页合并后，需要同步本地：
\begin{lstlisting}
git switch main
git -c http.proxy=http://127.0.0.1:7897 pull --ff-only origin main
git branch -d feature/section-area
\end{lstlisting}

% ============================================================
% 1.14 Jupyter Notebook 项目管理（第十一课）
% ============================================================
\section{Jupyter Notebook 与科研项目管理}
\subsection{Python 脚本与 Notebook 的区别}
\texttt{.py} 适合可复用的算法与程序，\texttt{.ipynb} 适合交互式计算、Markdown 说明、公式与可视化。Notebook 文件还可能保存运行输出和执行计数，因此每次运行后 Git 差异可能较大。
\subsection{科研仓库目录组织}
\begin{lstlisting}
my-python-learning/
  code/               # Python source code
  notebooks/          # Jupyter notebooks
  data/               # Small, shareable datasets
  results/            # Selected outputs
  README.md
  .gitignore
\end{lstlisting}
\subsection{Notebook 中的科研计算示例}
在 \texttt{notebooks/01\_python\_basics.ipynb} 中新建代码单元：
\begin{lstlisting}[language=Python]
import numpy as np
import matplotlib.pyplot as plt
length, width = 400, 50
thicknesses = np.array([0.3, 0.4, 0.5, 0.6, 0.7])
volumes = length * width * thicknesses
plt.plot(thicknesses, volumes, "o-")
plt.xlabel("Thickness (mm)")
plt.ylabel("Volume (mm^3)")
plt.grid(True)
plt.show()
\end{lstlisting}
\subsection{提交 Notebook 前的检查}
建议清除不必要的运行输出，检查文件是否包含本地路径、令牌或敏感数据。只添加需要上传的 Notebook，避免直接执行 \texttt{git add .} 误提交其他文件。
\begin{lstlisting}
git status
git add notebooks/01_python_basics.ipynb
git diff --staged --stat
git commit -m "Add first scientific computing notebook"
git -c http.proxy=http://127.0.0.1:7897 push origin main
\end{lstlisting}
\subsection{科研项目可复现性}
建议记录代码版本、数据来源、Python 环境与依赖库。大型数据和训练模型应根据大小与保密要求采用独立存储或 Git LFS，不应直接将涉密数据上传到 GitHub。

% ============================================================

% 1.15 本章总结

% ============================================================



\section{本章总结}



Git 的核心操作流程为：



\begin{center}

\texttt{git status}

$\longrightarrow$

\texttt{git add}

$\longrightarrow$

\texttt{git commit}

$\longrightarrow$

\texttt{git push}

\end{center}



本章介绍了 Git 的基础概念、

安装配置、仓库克隆、本地提交、

远程同步、版本恢复和文件忽略机制。



在后续 Python 与机器学习学习中，

可以利用 Git 管理算法实现、

实验脚本及配置文件，

并通过 GitHub 或 Gitee 保存和共享代码。



建议养成以下习惯：



\begin{enumerate}

    \item 修改代码前检查仓库状态；

    \item 提交前检查代码差异；

    \item 每次提交使用明确的说明；

    \item 定期同步远程仓库；

    \item 不上传密码、访问令牌及敏感数据。

\end{enumerate}
